\section{Terminology for readers new to the area}
\label{app:glossary}

This appendix defines every term the paper uses, assuming no background.
Nothing here is needed by a reader who already works on language models;
it exists so the main text can stay short.

\subsection{The problem}

\begin{description}[leftmargin=0em,style=nextline,itemsep=0.3em]

\item[State] The input being judged: a customer message, a product review,
a comment. One state can have many questions asked about it.

\item[Question] One decision to make about the state, with instructions
and a set of candidate answers.

\item[Option, menu, criteria] Used interchangeably for the candidate
answers of a question. A \emph{fixed menu} is shared by every row of a
task; a \emph{per-example menu} differs row by row, as in reading
comprehension where each question has its own answers.

\item[$K$] The number of options. $K{=}151$ means a 151-way choice.

\item[Primitive] The type of a question: \choice{} (pick one of $K$),
\score{} (pick one level on an ordered scale), \noul{} (yes or no).

\item[Multi-label] When more than one label can be true at once. A single
\choice{} cannot express this, because its probabilities sum to one across
options. Several \noul{} questions can.

\end{description}

\subsection{How models process text}

\begin{description}[leftmargin=0em,style=nextline,itemsep=0.3em]

\item[Token] A chunk of text, roughly a word piece. Models read and write
sequences of tokens, not characters.

\item[Embedding, hidden state] The vector a model holds for each token
position after processing. The hidden state at a position is what we read
a decision from.

\item[Attention] The mechanism by which a token's representation is built
from other tokens. An \emph{attention mask} says which positions may look
at which.

\item[Encoder vs decoder] Two attention patterns. An \textbf{encoder} is
bidirectional: every token sees every other. A \textbf{decoder} is causal:
a token sees only what came before it, which is what makes text generation
possible. \textbf{In this paper ``decoder'' names the attention pattern
only.} We never generate text; we run one forward pass and read scores.
The decoder is the backbone we recommend; the encoder is the documented
alternative (\S\ref{sec:when-backbone}).

\item[Forward pass] One run of the model over an input. Our central
efficiency claim is that $m$ questions cost one forward pass, not $m$.

\item[Logit] A raw, unnormalised score. Logits become probabilities
through a softmax.

\item[Softmax] Turns a vector of logits into probabilities that sum to
one. Our \emph{grouped} softmax does this separately within each
question's options, because questions in one packed sequence have
different $K$.

\item[Backbone] The pretrained model we build on, before any of our
training.

\end{description}

\subsection{Training}

\begin{description}[leftmargin=0em,style=nextline,itemsep=0.3em]

\item[Pretraining] The expensive first stage, done by others, where a
model learns language from large text corpora. We never do this.

\item[Fine-tuning] Further training of a pretrained model on a specific
task. Cheap by comparison: 38 seconds in one of our settings.

\item[Epoch] One pass over the training data.

\item[Learning rate] How large a step training takes per update. It
matters more than it sounds: our LoRA-beats-full-fine-tuning result is
most likely a learning-rate effect (\S\ref{sec:res-router}).

\item[Frozen] A parameter that is not updated during training.

\item[LoRA (Low-Rank Adaptation)] Instead of updating a large weight
matrix, train two small matrices alongside it and add their product. The
large matrix stays frozen. The \textbf{rank} $r$ sets how small; $r{=}16$
here. The result is an \textbf{adapter} of tens of megabytes rather than
gigabytes, so one backbone can serve many tasks.

\item[Merging] Folding a trained adapter back into the frozen weights, so
inference costs no extra work. An unmerged adapter is roughly twice as
slow for no benefit (\S\ref{sec:res-latency}).

\item[Multi-task mixture] Training on many tasks at once, so the model
learns the general shape of the problem rather than one instance of it.

\item[Intermediate-task transfer] Train on task A, then on target task B.
Helps most when B is small. Known in the literature as
STILTs~\cite{stilts}.

\item[Negative transfer] When adding training data makes the target
\emph{worse}. We caused this twice (\S\ref{sec:ablations}).

\end{description}

\subsection{Evaluation}

\begin{description}[leftmargin=0em,style=nextline,itemsep=0.3em]

\item[Zero-shot] Evaluating on a task the model was never trained on. The
hard claim, and the thing Jev is good at.

\item[Held-out] Data deliberately excluded from training, used to measure
zero-shot ability. Worthless if it leaks into training, hence the
contamination guard.

\item[Held-in control] Accuracy on tasks the model \emph{was} trained on,
reported next to held-out. If held-out is at chance and held-in is high,
the model truly fails to transfer; if both are at chance, something is
broken. Without this, a bug and a finding look identical.

\item[Chance] What random guessing scores: $1/K$. On a 151-way menu that
is $0.007$, so an accuracy of $0.29$ is strong; on a 4-way menu chance is
$0.25$ and $0.29$ is nearly worthless.

\item[Multiple of chance] Accuracy divided by chance, so tasks with very
different $K$ can be compared on one axis.

\item[Majority-class baseline] Always predicting the most common label.
On skewed data this beats chance, so it is the stricter floor. One of our
encoder results passes the chance test and only ties this one
(\S\ref{sec:abl-encoder}).

\item[Accuracy, precision, recall] Accuracy is the fraction correct.
\emph{Recall} is the fraction of true positives found, which is what
matters for a safety gate. \emph{Precision} is the fraction of positive
predictions that were right.

\item[Macro-F1] The harmonic mean of precision and recall, averaged over
classes so rare classes count as much as common ones. A binary task
scoring macro-F1 $0.333$ is predicting one class for everything.

\item[AUROC] Ranking quality, independent of any threshold: the
probability a random positive is scored above a random negative. $0.5$ is
chance. It matters because a model can rank well and still score badly at
a fixed threshold, which is a calibration problem and not a capability
one.

\item[Calibration] Whether a stated probability matches reality. If
everything a model calls $0.7$ is right $70\%$ of the time, it is
calibrated. \textbf{ECE} measures the gap; \textbf{temperature scaling} is
the standard post-hoc fix.

\item[Bootstrap confidence interval] Resample the test set with
replacement many times and report the middle 95\% of the resulting scores.
An interval that excludes chance means the result is unlikely to be luck.

\item[McNemar test] For comparing two systems on the \emph{same} items.
Counts only the disagreements: $b_{10}$ where we are right and they are
not, $b_{01}$ the reverse, and asks whether the split is lopsided enough
to be real. Necessary because two systems can have equal accuracy while
disagreeing on a third of the set.

\item[$p$-value] Informally, the probability of seeing a difference this
large if there were no real difference. Below $0.05$ is conventionally
called significant. With small samples, significance can be unreachable
no matter how good the model: one of our tiers has 29 items, so even a
perfect score floors at $p = 0.0625$.

\end{description}

\subsection{Terms specific to this work}

\begin{description}[leftmargin=0em,style=nextline,itemsep=0.3em]

\item[Packing] Writing the state once, then all questions and options
after it, as one sequence for one forward pass.

\item[Shared prefix] The state portion of a packed sequence, attended to
by every question. The reason $m$ questions are nearly free.

\item[Marker] A token after each option where its score is read.

\item[Block-diagonal attention] The mask that lets each question see the
shared prefix and itself, but not other questions, so an answer does not
depend on what else was asked.

\item[Label-space augmentation] Padding a training menu with labels
borrowed from other tasks. The gold stays correct, but the model must
discriminate against a larger menu. Our largest single effect.

\item[Distractor] One of those borrowed labels.

\item[Contamination guard] A check that no evaluation dataset appears in
the training mixture under any of its names.

\end{description}
